\section{Tarefa 1 --- Modelagem com MER/DER}

\subsection{Entidades e atributos}

A Tabela~\ref{tab:entidades} apresenta as entidades do modelo, seus atributos e a chave primária. As chaves estrangeiras estão indicadas por ``(FK)'' e o atributo derivado por ``(derivado)''. Para manter o DER legível, os atributos são detalhados na tabela e não repetidos na visão geral da Figura~\ref{fig:der}.

{\footnotesize
\begin{longtable}{@{}L{3.5cm}L{9.4cm}L{1.9cm}@{}}
\caption{Entidades, atributos e chaves}
\label{tab:entidades}\\
\toprule
\textbf{Entidade} & \textbf{Atributos} & \textbf{Chave primária}\\
\midrule
\endfirsthead
\toprule
\textbf{Entidade} & \textbf{Atributos} & \textbf{Chave primária}\\
\midrule
\endhead
\bottomrule
\endfoot

\texttt{Cliente} &
\texttt{cpf}; \texttt{nome\_completo}; \texttt{telefone}; \texttt{email}; \texttt{endereco} (logradouro, número, complemento, bairro, cidade, UF, CEP). &
\texttt{cpf}\\

\texttt{Coleta} &
\texttt{codigo\_coleta}; \texttt{data\_coleta}; \texttt{data\_entrega\_coleta}; \texttt{cpf\_remetente} (FK); \texttt{cpf\_destinatario} (FK); \texttt{nome\_remetente}; \texttt{endereco\_coleta}; \texttt{telefone\_remetente}; \texttt{nome\_destinatario}; \texttt{endereco\_entrega}; \texttt{telefone\_destinatario}; \texttt{codigo\_deposito\_origem} (FK); \texttt{codigo\_cidade\_destino} (FK). &
\texttt{codigo\_coleta}\\

\texttt{Encomenda} &
\texttt{codigo\_encomenda}; \texttt{codigo\_coleta} (FK); \texttt{descricao}; \texttt{valor}; \texttt{peso\_kg}; \texttt{volume\_m3}; \texttt{tipo\_encomenda} (discriminador). &
\texttt{codigo\_encomenda}\\

\texttt{Encomenda\allowbreak Comum} &
\texttt{codigo\_encomenda} (PK/FK). Sem atributos específicos. &
\texttt{codigo\_encomenda}\\

\texttt{Encomenda\allowbreak Fragil} &
\texttt{codigo\_encomenda} (PK/FK); \texttt{cuidados\_especiais}; \texttt{empilhamento\_maximo}; \texttt{manter\_em\_pe}. &
\texttt{codigo\_encomenda}\\

\texttt{Encomenda\allowbreak Perecivel} &
\texttt{codigo\_encomenda} (PK/FK); \texttt{temperatura\_min\_c}; \texttt{temperatura\_max\_c}; \texttt{validade}. &
\texttt{codigo\_encomenda}\\

\texttt{Encomenda\allowbreak Perigosa} &
\texttt{codigo\_encomenda} (PK/FK); \texttt{numero\_onu}; \texttt{classe\_risco}; \texttt{grupo\_embalagem}. &
\texttt{codigo\_encomenda}\\

\texttt{NotaFiscal} &
\texttt{numero\_nota}; \texttt{codigo\_encomenda} (FK, único); \texttt{codigo\_fatura} (FK); \texttt{data\_emissao}; \texttt{valor\_total}; \texttt{percentual\_frete}; \texttt{valor\_frete} (derivado); \texttt{data\_entrega\_destinatario}. &
\texttt{numero\_nota}\\

\texttt{Fatura} &
\texttt{numero\_fatura}; \texttt{cpf\_pagador} (FK); \texttt{data\_emissao}; \texttt{data\_vencimento}; \texttt{valor\_frete\_total} (derivado). &
\texttt{numero\_fatura}\\

\texttt{Caminhao} &
\texttt{placa}; \texttt{modelo}; \texttt{capacidade\_kg}; \texttt{capacidade\_m3}. &
\texttt{placa}\\

\texttt{Cidade} &
\texttt{codigo\_cidade}; \texttt{nome}; \texttt{uf}; \texttt{codigo\_regiao} (FK). &
\texttt{codigo\_cidade}\\

\texttt{Regiao} &
\texttt{codigo\_regiao}; \texttt{nome}. &
\texttt{codigo\_regiao}\\

\texttt{Deposito} &
\texttt{codigo\_deposito}; \texttt{nome}; \texttt{endereco}; \texttt{codigo\_cidade} (FK, único). &
\texttt{codigo\_deposito}\\

\texttt{Embarque} &
\texttt{codigo\_embarque}; \texttt{numero\_nota} (FK); \texttt{placa} (FK); \texttt{codigo\_cidade\_destino} (FK); \texttt{ordem\_trecho}; \texttt{data\_hora\_embarque}; \texttt{data\_hora\_desembarque}. &
\texttt{codigo\_embarque}\\

\texttt{Ocorrencia} &
\texttt{codigo\_ocorrencia}; \texttt{codigo\_embarque} (FK); \texttt{tipo}; \texttt{data\_hora}; \texttt{descricao}; \texttt{status}. &
\texttt{codigo\_ocorrencia}\\
\end{longtable}
}

\noindent\footnotesize Os campos \texttt{nome\_remetente}, \texttt{endereco\_coleta}, \texttt{telefone\_remetente}, \texttt{nome\_destinatario}, \texttt{endereco\_entrega} e \texttt{telefone\_destinatario} são cópias históricas dos dados do cliente no momento da coleta, exigidas pelo documento de coleta. \normalsize

\subsection{Relacionamentos e cardinalidades}

A Tabela~\ref{tab:relacionamentos} resume os relacionamentos, as cardinalidades e a participação. A notação $(m,n)$ indica o mínimo e o máximo de ocorrências de uma entidade associadas a uma ocorrência da outra.

{\small
\begin{longtable}{@{}L{3.0cm}L{3.0cm}L{2.7cm}L{6.2cm}@{}}
\caption{Relacionamentos, cardinalidades e participação}
\label{tab:relacionamentos}\\
\toprule
\textbf{Relacionamento} & \textbf{Entidades} & \textbf{Cardinalidade} & \textbf{Regra de negócio (participação)}\\
\midrule
\endfirsthead
\toprule
\textbf{Relacionamento} & \textbf{Entidades} & \textbf{Cardinalidade} & \textbf{Regra de negócio (participação)}\\
\midrule
\endhead
\bottomrule
\endfoot

\texttt{ENVIA} & Cliente--Coleta & 1:N & Total em Coleta: um cliente pode ser remetente de várias coletas; cada coleta tem um remetente.\\

\texttt{RECEBE} & Cliente--Coleta & 1:N & Total em Coleta: um cliente pode ser destinatário de várias coletas; cada coleta tem um destinatário.\\

\texttt{CONTEM} & Coleta--Encomenda & 1:N & Total em Encomenda: uma coleta contém uma ou mais encomendas; cada encomenda pertence a uma coleta.\\

\texttt{GERA} & Encomenda--NotaFiscal & 1:1 & Total nos dois lados: cada encomenda gera uma nota; cada nota se refere a uma encomenda.\\

\texttt{COBRA} & NotaFiscal--Fatura & N:1 & Total em NotaFiscal: uma fatura agrupa várias notas; cada nota pertence a uma fatura.\\

\texttt{PAGA} & Cliente--Fatura & 1:N & Total em Fatura: o pagador é o remetente ou o destinatário; um cliente pode pagar várias faturas.\\

\texttt{ORIGINA} & Deposito--Coleta & 1:N & Total em Coleta: cada coleta é entregue a um depósito; um depósito recebe várias coletas.\\

\texttt{ENTREGA\_EM} & Coleta--Cidade & N:1 & Total em Coleta: cada coleta tem uma cidade de destino; uma cidade recebe várias coletas.\\

\texttt{SEDIA} & Deposito--Cidade & 1:1 & Total em Deposito e parcial em Cidade: no máximo um depósito por cidade; todo depósito fica em uma cidade.\\

\texttt{AGRUPIA} & Cidade--Regiao & N:1 & Total em Cidade: cada cidade pertence a uma região; uma região agrupa várias cidades.\\

\texttt{ATENDE} & Caminhao--Cidade & N:N & Total em Caminhao, parcial em Cidade: um caminhão atende uma ou mais cidades (P6); uma cidade pode não ter caminhões cadastrados.\\

\texttt{EMBARCA} & NotaFiscal--Caminhao & N:N & Total em NotaFiscal, parcial em Caminhao: uma nota tem um ou mais embarques (P7); um caminhão pode não ter embarques.\\

\texttt{DESTINO\_TRECHO} & Embarque--Cidade & N:1 & Total em Embarque: cada trecho tem uma cidade de destino.\\

\texttt{REGISTRA} & Embarque--Ocorrencia & 1:N & Total em Ocorrencia: um embarque pode ter várias ocorrências; cada ocorrência pertence a um embarque.\\
\end{longtable}
}

\subsection{Diagramas}

A Figura~\ref{fig:der} apresenta a visão geral do DER. Para evitar poluição visual, os atributos foram omitidos dessa figura; eles estão na Tabela~\ref{tab:entidades} e nos detalhes das Figuras~\ref{fig:hierarquia} e~\ref{fig:agregacao}. Todas as entidades, relacionamentos e cardinalidades do modelo estão representados.

\begin{figure}[H]
\centering
\resizebox{\linewidth}{!}{\begin{tikzpicture}[x=1cm,y=1cm,every node/.style={transform shape}]
  % ---------- Entidades ----------
  \node[ent] (cliente) at (-10.2,5.0) {Cliente};
  \node[ent] (coleta) at (-5.0,5.0) {Coleta};
  \node[ent] (encomenda) at (-0.2,5.0) {Encomenda};
  \node[ent] (nota) at (4.6,5.0) {NotaFiscal};
  \node[ent] (fatura) at (9.8,5.0) {Fatura};

  \node[entref] (deposito) at (-10.2,1.0) {Deposito};
  \node[entref] (cidade) at (-2.2,1.0) {Cidade};
  \node[entref] (regiao) at (4.2,1.0) {Regiao};

  \node[ent] (caminhao) at (-10.0,-3.6) {Caminhao};
  \node[entadd] (embarque) at (0.6,-3.6) {Embarque};
  \node[entadd] (ocorrencia) at (7.6,-3.6) {Ocorrencia};

  % ---------- Relacionamentos ----------
  \node[rel] (envia) at (-7.6,6.25) {ENVIA};
  \node[rel] (recebe) at (-7.6,3.75) {RECEBE};
  \node[rel] (origina) at (-8.2,2.05) {ORIGINA};
  \node[rel] (contem) at (-2.6,5.0) {CONTEM};
  \node[rel] (gera) at (2.2,5.0) {GERA};
  \node[rel] (cobra) at (7.2,5.0) {COBRA};
  \node[rel] (entrega) at (-3.7,3.0) {ENTREGA\\EM};
  \node[rel] (sedia) at (-6.2,1.0) {SEDIA};
  \node[rel] (agrupia) at (1.0,1.0) {AGRUPIA};
  \node[rel] (atende) at (-6.2,-1.3) {ATENDE};
  \node[reladd] (embarca) at (-2.8,-2.8) {EMBARCA};
  \node[rel] (destino) at (-0.4,-1.3) {DESTINO\\TRECHO};
  \node[rel] (registra) at (4.0,-3.6) {REGISTRA};
  \node[rel] (paga) at (-4.5,7.3) {PAGA};

  % ---------- Ligações ----------
  \draw[linha] (cliente.east) -- (envia.west);
  \draw[linha] (envia.east) -- (coleta.west);
  \draw[linha] (cliente.east) -- (recebe.west);
  \draw[linha] (recebe.east) -- (coleta.west);

  \draw[linha] (deposito.north) -- (origina.west);
  \draw[linha] (origina.east) -- (coleta.south west);

  \draw[linha] (coleta.east) -- (contem.west);
  \draw[linha] (contem.east) -- (encomenda.west);
  \draw[linha] (encomenda.east) -- (gera.west);
  \draw[linha] (gera.east) -- (nota.west);
  \draw[linha] (nota.east) -- (cobra.west);
  \draw[linha] (cobra.east) -- (fatura.west);

  \draw[linha] (coleta.south) -- (entrega.north);
  \draw[linha] (entrega.south) -- (cidade.north west);

  \draw[linha] (deposito.east) -- (sedia.west);
  \draw[linha] (sedia.east) -- (cidade.west);
  \draw[linha] (cidade.east) -- (agrupia.west);
  \draw[linha] (agrupia.east) -- (regiao.west);

  \draw[linha] (caminhao.north) -- (atende.south);
  \draw[linha] (atende.north) -- (cidade.south west);

  % EMBARCA como agregação entre NotaFiscal e Caminhao
  \draw[linha] (nota.south west) -- (embarca.north east);
  \draw[linha] (caminhao.north east) -- (embarca.south west);
  \draw[linha] (embarca.east) -- (embarque.west);

  \draw[linha] (embarque.north) -- (destino.south);
  \draw[linha] (destino.north) -- (cidade.south east);

  \draw[linha] (embarque.east) -- (registra.west);
  \draw[linha] (registra.east) -- (ocorrencia.west);

  % Pagador: rota superior
  \draw[linha] (fatura.north) -- (9.8,7.3) -- (paga.east);
  \draw[linha] (paga.west) -- (-10.2,7.3) -- (cliente.north);

  % ---------- Cardinalidades ----------
  \node[card, above left] at (envia.west) {(0,N)};
  \node[card, above right] at (envia.east) {(1,1)};
  \node[card, below left] at (recebe.west) {(0,N)};
  \node[card, below right] at (recebe.east) {(1,1)};
  \node[card, left] at (origina.west) {(0,N)};
  \node[card, right] at (origina.east) {(1,1)};
  \node[card] at (-3.35,5.62) {(1,N)};
  \node[card] at (-1.85,5.62) {(1,1)};
  \node[card] at (1.6,5.62) {(1,1)};
  \node[card] at (2.8,5.62) {(1,1)};
  \node[card] at (6.5,5.62) {(1,N)};
  \node[card] at (7.9,5.62) {(1,1)};
  \node[card, left] at (entrega.north) {(1,1)};
  \node[card, right] at (entrega.south) {(0,N)};
  \node[card, below] at (sedia.west) {(0,1)};
  \node[card, below] at (sedia.east) {(1,1)};
  \node[card, below] at (agrupia.west) {(0,N)};
  \node[card, below] at (agrupia.east) {(1,1)};
  \node[card, left] at (atende.south) {(1,N)};
  \node[card, right] at (atende.north) {(0,N)};
  % EMBARCA cardinalidades
  \node[card, above left] at (embarca.north west) {(1,N)};
  \node[card, below right] at (embarca.south east) {(0,N)};
  \node[card, left] at (destino.south) {(1,1)};
  \node[card, right] at (destino.north) {(0,N)};
  \node[card, above] at (3.3,-3.15) {(0,N)};
  \node[card, above] at (4.75,-3.15) {(1,1)};
  \node[card, above left] at (paga.west) {(0,N)};
  \node[card, above right] at (paga.east) {(1,1)};

  % ---------- Papéis e anotações ----------
  \node[rotulo, anchor=west] at (fatura.east) {pagador};
  \node[agregacao, fit=(embarca), inner sep=0.18cm] {};

  % ---------- Legenda ----------
  \node[draw=cinza, rounded corners=2pt, fill=white, text width=7.2cm,
        align=left, font=\scriptsize, inner sep=3pt] at (-7.6,-6.35) {
    \textbf{Legenda.} Retângulo: entidade. Losango: relacionamento.
    Linha dupla: participação total. $(m,n)$: cardinalidade mínima e máxima.
    Caixa tracejada: agregação. Cores: base (azul), referência (verde),
    adicional (laranja).
  };
\end{tikzpicture}}
\caption{Visão geral do DER do Sistema de Transportadora}
\label{fig:der}
\end{figure}

A Figura~\ref{fig:hierarquia} detalha o requisito adicional RA1. A especialização é representada pelo círculo com a letra ``d'', que indica especialização disjunta; a linha dupla entre o supertipo e o círculo indica participação total. A Figura~\ref{fig:agregacao} detalha o requisito adicional RA2: o losango \texttt{EMBARCA} é envolvido pela caixa tracejada (agregação) e materializado pela entidade \texttt{Embarque}, que por sua vez se relaciona com \texttt{Ocorrencia}.

\begin{figure}[H]
\centering
\resizebox{0.82\linewidth}{!}{\begin{tikzpicture}[x=1cm,y=1cm,every node/.style={transform shape}]
  % Supertipo
  \node[draw=azul, thick, rounded corners=2pt, fill=azulclaro,
        text width=8.4cm, align=left, font=\scriptsize, inner sep=4pt]
        (enc) at (0,4.9) {
    \textbf{\normalsize Encomenda}\\[-2pt]
    \rule{8.0cm}{0.4pt}\\[-3pt]
    \underline{\texttt{codigo\_encomenda}} (PK)\\
    \texttt{codigo\_coleta} (FK) \quad \texttt{descricao}\\
    \texttt{valor} \quad \texttt{peso\_kg} \quad \texttt{volume\_m3}\\
    \texttt{tipo\_encomenda} (discriminador)
  };

  % Circulo de especializacao
  \node[draw=roxo, thick, circle, fill=roxoclaro, minimum size=0.75cm,
        font=\small\bfseries] (esp) at (0,3.05) {d};
  \draw[draw=cinza, thick, double, double distance=1.3pt] (enc.south) -- (esp.north);
  \node[font=\scriptsize\itshape, anchor=west] at (0.45,3.05) {total, disjunta};

  % Barra de distribuicao
  \draw[draw=cinza, thick] (esp.south) -- (0,2.1);
  \draw[draw=cinza, thick] (-6.3,2.1) -- (6.3,2.1);

  % Subtipos
  \node[entsub, text width=3.75cm] (comum) at (-6.3,0.15) {
    \textbf{\footnotesize EncomendaComum}\\[-1pt]
    \rule{3.5cm}{0.3pt}\\[-2pt]
    \underline{\texttt{codigo\_encomenda}} (PK/FK)\\
    \textit{Sem atributos específicos}
  };
  \node[entsub, text width=3.75cm] (fragil) at (-2.1,0.15) {
    \textbf{\footnotesize EncomendaFragil}\\[-1pt]
    \rule{3.5cm}{0.3pt}\\[-2pt]
    \underline{\texttt{codigo\_encomenda}} (PK/FK)\\
    \texttt{cuidados\_especiais}\\
    \texttt{empilhamento\_maximo}\\
    \texttt{manter\_em\_pe}
  };
  \node[entsub, text width=3.75cm] (perec) at (2.1,0.15) {
    \textbf{\footnotesize EncomendaPerecivel}\\[-1pt]
    \rule{3.5cm}{0.3pt}\\[-2pt]
    \underline{\texttt{codigo\_encomenda}} (PK/FK)\\
    \texttt{temperatura\_min\_c}\\
    \texttt{temperatura\_max\_c}\\
    \texttt{validade}
  };
  \node[entsub, text width=3.75cm] (perig) at (6.3,0.15) {
    \textbf{\footnotesize EncomendaPerigosa}\\[-1pt]
    \rule{3.5cm}{0.3pt}\\[-2pt]
    \underline{\texttt{codigo\_encomenda}} (PK/FK)\\
    \texttt{numero\_onu}\\
    \texttt{classe\_risco}\\
    \texttt{grupo\_embalagem}
  };

  % Liga barra aos subtipos
  \draw[draw=cinza, thick] (-6.3,2.1) -- (comum.north);
  \draw[draw=cinza, thick] (-2.1,2.1) -- (fragil.north);
  \draw[draw=cinza, thick] (2.1,2.1) -- (perec.north);
  \draw[draw=cinza, thick] (6.3,2.1) -- (perig.north);

  % Nota
  \node[draw=cinza, dashed, rounded corners=2pt, fill=white, text width=13.6cm,
        align=center, font=\scriptsize, inner sep=3pt] at (0,-2.2) {
    Toda encomenda pertence a exatamente um dos quatro subtipos (especialização total e disjunta).
    O discriminador \texttt{tipo\_encomenda} deve corresponder ao subtipo instanciado.
  };
\end{tikzpicture}
}
\caption{Hierarquia de especialização de \texttt{Encomenda} (RA1)}
\label{fig:hierarquia}
\end{figure}

\begin{figure}[H]
\centering
\resizebox{0.78\linewidth}{!}{\begin{tikzpicture}[x=1cm,y=1cm]
  % Entidades participantes
  \node[draw=azul, thick, rounded corners=2pt, fill=azulclaro,
        text width=4.0cm, align=left, font=\scriptsize, inner sep=3pt]
        (nota) at (-5.3,5.6) {
    \textbf{\normalsize NotaFiscal}\\[-2pt]
    \rule{3.7cm}{0.4pt}\\[-3pt]
    \underline{\texttt{numero\_nota}} (PK)\\
    \texttt{data\_emissao}\\
    \texttt{valor\_total}\\
    \texttt{valor\_frete}\\
    \texttt{data\_entrega\_destinatario}
  };
  \node[draw=azul, thick, rounded corners=2pt, fill=azulclaro,
        text width=3.4cm, align=left, font=\scriptsize, inner sep=3pt]
        (cam) at (5.3,5.6) {
    \textbf{\normalsize Caminhao}\\[-2pt]
    \rule{3.1cm}{0.4pt}\\[-3pt]
    \underline{\texttt{placa}} (PK)\\
    \texttt{modelo}\\
    \texttt{capacidade\_kg}\\
    \texttt{capacidade\_m3}
  };

  % Relacionamento agregado
  \node[reladd, minimum width=1.9cm, minimum height=0.85cm] (embarca) at (0,5.6) {EMBARCA};
  \draw[draw=cinza, thick] (nota.east) -- (embarca.west);
  \draw[draw=cinza, thick] (embarca.east) -- (cam.west);
  \node[card, above=2pt of embarca] {N:N};
  \node[card, left=3pt of embarca] {(0,N)};
  \node[card, right=3pt of embarca] {(0,N)};

  % Caixa de agregacao
  \node[agregacao, fit=(embarca), inner sep=0.55cm] (box) {};
  \node[font=\scriptsize\itshape, text=laranja, anchor=south] at (box.north) {agregação};

  % Entidade associativa materializada
  \node[draw=laranja, thick, rounded corners=2pt, fill=laranjaclaro,
        text width=5.2cm, align=left, font=\scriptsize, inner sep=3pt]
        (emb) at (0,2.4) {
    \textbf{\normalsize Embarque}\\[-2pt]
    \rule{4.9cm}{0.4pt}\\[-3pt]
    \underline{\texttt{codigo\_embarque}} (PK)\\
    \texttt{numero\_nota} (FK) \quad \texttt{placa} (FK)\\
    \texttt{codigo\_cidade\_destino} (FK)\\
    \texttt{ordem\_trecho}\\
    \texttt{data\_hora\_embarque}\\
    \texttt{data\_hora\_desembarque}
  };
  \draw[seta] (embarca.south) -- (emb.north);
  \node[rotulo, anchor=west] at (1.55,4.95) {materializa};

  % Ocorrencia
  \node[rel, minimum width=1.7cm, minimum height=0.75cm] (registra) at (0,-0.4) {REGISTRA};
  \node[draw=laranja, thick, rounded corners=2pt, fill=laranjaclaro,
        text width=4.8cm, align=left, font=\scriptsize, inner sep=3pt]
        (oco) at (0,-2.5) {
    \textbf{\normalsize Ocorrencia}\\[-2pt]
    \rule{4.5cm}{0.4pt}\\[-3pt]
    \underline{\texttt{codigo\_ocorrencia}} (PK)\\
    \texttt{codigo\_embarque} (FK)\\
    \texttt{tipo} \quad \texttt{data\_hora}\\
    \texttt{descricao} \quad \texttt{status}
  };
  \draw[draw=cinza, thick, double, double distance=1.2pt] (emb.south) -- (registra.north);
  \draw[draw=cinza, thick, double, double distance=1.2pt] (registra.south) -- (oco.north);
  \node[card, anchor=west] at (0.6,0.05) {(0,N)};
  \node[card, anchor=west] at (0.6,-0.95) {(1,1)};
\end{tikzpicture}
}
\caption{Agregação \texttt{EMBARCA} e entidades \texttt{Embarque} e \texttt{Ocorrencia} (RA2)}
\label{fig:agregacao}
\end{figure}

\subsection{Como os requisitos da Seção 3 foram modelados}

A Tabela~\ref{tab:requisitos_mer} relaciona cada item do enunciado à decisão de modelagem correspondente. Os requisitos adicionais RA1 e RA2 estão descritos na Seção~\ref{sec:adicionais} e representados nas Figuras~\ref{fig:hierarquia} e~\ref{fig:agregacao}.

{\small
\begin{longtable}{@{}L{1.0cm}L{4.6cm}L{9.4cm}@{}}
\caption{Correspondência entre os requisitos da Seção 3 e o MER}
\label{tab:requisitos_mer}\\
\toprule
\textbf{Item} & \textbf{Requisito} & \textbf{Como foi modelado no MER}\\
\midrule
\endfirsthead
\toprule
\textbf{Item} & \textbf{Requisito} & \textbf{Como foi modelado no MER}\\
\midrule
\endhead
\bottomrule
\endfoot

R01 & O cliente contata a transportadora e é registrada uma coleta no endereço indicado. & Entidade \texttt{Coleta}, com \texttt{data\_coleta} e \texttt{endereco\_coleta}, relacionada ao cliente remetente por \texttt{ENVIA}.\\

R02 & O cliente é o remetente e o destinatário também é um cliente. & Dois relacionamentos 1:N entre \texttt{Cliente} e \texttt{Coleta}: \texttt{ENVIA} (remetente) e \texttt{RECEBE} (destinatário).\\

R03 & Guardar nome completo, CPF, telefone, e-mail e endereço completo do remetente e do destinatário. & Atributos em \texttt{Cliente}, com \texttt{cpf} como chave primária. O endereço é um atributo composto.\\

R04 & O documento de coleta registra CPF, nome, endereço e telefone de remetente e destinatário. & \texttt{Coleta} guarda as FKs \texttt{cpf\_remetente} e \texttt{cpf\_destinatario} e cópias históricas de nome, endereço e telefone de ambos.\\

R05 & Dados da mercadoria: descrição, valor, peso e volume em metros cúbicos. & Entidade \texttt{Encomenda}, com \texttt{descricao}, \texttt{valor}, \texttt{peso\_kg} e \texttt{volume\_m3}, em relacionamento 1:N com \texttt{Coleta} (\texttt{CONTEM}).\\

R06 & Emitir nota fiscal com os dados do remetente e do destinatário, buscando os dados da coleta, inclusive a mercadoria. & \texttt{NotaFiscal} em 1:1 com \texttt{Encomenda} (\texttt{GERA}); os dados de remetente e destinatário são obtidos pelo caminho \texttt{NotaFiscal}$\rightarrow$\texttt{Encomenda}$\rightarrow$\texttt{Coleta}$\rightarrow$\texttt{Cliente}, sem duplicação no MER.\\

R07 & Os caminhões que fazem coletas em uma cidade chegam ao depósito; há um depósito por cidade; as mercadorias são agrupadas por cidade ou região de destino. & \texttt{Deposito} em 1:1 com \texttt{Cidade} (\texttt{SEDIA}); \texttt{Coleta} se relaciona com \texttt{Deposito} (\texttt{ORIGINA}) e com a cidade de destino (\texttt{ENTREGA\_EM}); \texttt{Cidade} se relaciona com \texttt{Regiao} (\texttt{AGRUPIA}).\\

R08 & Para cada mercadoria é gerada uma nota fiscal. & Relacionamento 1:1 \texttt{GERA}, com participação total nos dois lados.\\

R09 & A coleta tem uma data de entrega. & Atributo \texttt{data\_entrega\_coleta} em \texttt{Coleta} (entrega ao depósito, conforme a premissa P9).\\

R10 & As mercadorias são carregadas em caminhões; cada caminhão vai para uma ou muitas cidades; mais de um caminhão pode ir para a mesma cidade. & Relacionamento N:N \texttt{ATENDE} entre \texttt{Caminhao} e \texttt{Cidade}; o transporte efetivo é registrado por \texttt{EMBARCA}, entre \texttt{NotaFiscal} e \texttt{Caminhao}, materializado por \texttt{Embarque}.\\

R11 & Calcular o valor do frete, normalmente 5\% do valor total da nota; a nota tem data de entrega ao destinatário. & Atributos \texttt{valor\_total}, \texttt{percentual\_frete}, \texttt{valor\_frete} (derivado) e \texttt{data\_entrega\_destinatario} em \texttt{NotaFiscal}.\\

R12 & A nota não é paga pelo cliente; ao gerar a nota é gerada uma fatura do frete, com número, datas de emissão e vencimento, valor e nome do pagador; a fatura pode envolver várias notas. & \texttt{Fatura} em 1:N com \texttt{NotaFiscal} (\texttt{COBRA}); \texttt{Cliente} se relaciona com \texttt{Fatura} no papel de pagador (\texttt{PAGA}); \texttt{valor\_frete\_total} é derivado das notas.\\
\end{longtable}
}
